
\section{Classical Mechanics as a Subject in Its Own Right}
\label{cm:st}

It is worth pausing on the first role only long enough to insist that it is real.
The three-body problem, the precession of orbits, the stability of the solar
system, the dynamics of rigid bodies and gyroscopes, the propagation of waves in
elastic media, the onset of turbulence, and the sensitive dependence on initial
conditions studied by dynamical-systems theory are all substantial, still
partially open, areas of physics conducted entirely within the classical
framework. None of this material is a mere warm-up exercise for quantum mechanics,
and a physicist could spend a full career inside it without ever needing $\hbar$.
This book is not about that role. It is mentioned here only to make clear that
when the rest of this book treats classical mechanics as a base to be generalized,
it is not thereby claiming that this exhausts what classical mechanics is.

The four cases below are developed fully in later chapters of this book. They
are introduced here, briefly and together, because their common moral is best
seen by contrast: in each case, a feature usually advertised as a signature of
quantum theory or of relativity turns out to have an exact, or near-exact,
analogue that requires nothing beyond Newton, Lagrange, and Hamilton.

\subsection{Hertz's forceless mechanics }

In his \emph{Principles of Mechanics}, Heinrich Hertz reformulated dynamics
without the concept of force at all. Observable bodies are treated as coupled to
unobservable ``hidden masses'' through rigid, holonomic constraints, and the
combined system moves along a geodesic of a metric built from the mass
distribution on an enlarged configuration space. Suitably arranged, Hertz's
constrained geodesic motion reproduces the trajectories that Newtonian mechanics
attributes to the action of a force. The two formulations are empirically
equivalent: no experiment performed on the observable bodies alone can decide
between ``a force is acting'' and ``a constraint is enforcing geodesic motion in
a larger space.''

This is a case of underdetermination of theory --- indeed of ontology --- by
observation, arising entirely within classical mechanics. It is normally quantum
mechanics that is credited with forcing physicists to confront multiple,
empirically indistinguishable accounts of what is ``really'' happening beneath
the observed statistics: Bohmian trajectories, Everettian branches, and
epistemic and objective-collapse readings of the wavefunction are all consistent
with the same experimental record. Hertz shows that the same structural
situation --- one set of observations, several inequivalent stories about the
underlying ontology --- already existed in nineteenth-century particle mechanics,
with no probability amplitude in sight.

\subsection{Nonholonomic constraints }

A constraint of the Pfaffian form
\begin{equation}
  \sum_i a_i(q)\,\dot q_i = 0
\end{equation}
is \emph{nonholonomic} when it is not integrable to a relation $f(q)=0$ among the
coordinates alone --- the standard example being a disk or coin rolling without
slipping on a plane. Such systems present a genuine fork in the road. Applying
d'Alembert's principle directly, with the constraint enforced at the level of
virtual displacements, yields the correct equations of motion, confirmed by
experiment. Applying Hamilton's variational principle instead --- extremizing the
action with the constraint imposed by a Lagrange multiplier, as one legitimately
may for holonomic constraints --- yields a different, and physically incorrect,
set of equations. This is the long-standing nonholonomic-versus-vakonomic
controversy: two superficially reasonable derivations from the same starting
Lagrangian and the same constraint disagree about the dynamics, and only one of
them is right.

One consequence is that nonholonomic systems generically fail to admit a
straightforward canonical Hamiltonian formulation on the naive phase space:
the constrained bracket structure one is forced to use is not the canonical
Poisson bracket, and which reduced structure is the correct one to use is itself
a nontrivial derivation, system by system. This matters directly for the
generalization picture of Section~\ref{cm:ba}: canonical quantization
needs a Poisson bracket to promote to a commutator, and here the classical
bracket is already ambiguous before any promotion is attempted. The
operator-ordering problem --- the fact that a classical observable such as $qp$
does not determine a unique quantum operator, since $\hat q\hat p$, $\hat p\hat
q$, and $\tfrac12(\hat q\hat p+\hat p\hat q)$ all reduce to it as $\hbar\to0$
--- is usually presented as an intrinsically quantum indeterminacy. In the
nonholonomic case the indeterminacy has already begun classically, in the
disagreement over what the correct pre-quantization bracket even is.

\subsection{Shape-changing bodies and classical geometric phase}

A body that changes its own shape --- the paradigm case is a cat, initially
oriented feet-up, that reorients itself in mid-air with zero net angular
momentum by cyclically changing its shape --- cannot be treated by rigid-body
mechanics. One instead works on a \emph{shape space} $\mathcal S$ (the space of
internal configurations, with rigid motions quotiented out) and finds that the
conservation of angular momentum, together with the equations of motion,
defines a connection on a principal bundle over $\mathcal S$ with structure
group $SO(3)$. A closed loop in shape space --- the cat returns to its initial
shape --- need not be a closed loop in the full configuration space: the body
can acquire a net rotation purely as a holonomy of this connection, with no
external torque and no change in angular momentum at any instant. The same
mechanism governs low-Reynolds-number swimmers and, more generally, any system
in which a periodic change of internal, ``slow'' variables produces a net shift
in an external, cyclic variable; the resulting angle is called the Hannay angle.

The geometric or Berry phase acquired by a quantum system whose Hamiltonian is
slowly and cyclically varied is frequently introduced as a paradigmatically
quantum-mechanical effect, precisely because it has no counterpart in the naive
classical picture of a particle following a trajectory. The falling cat and the
Hannay angle show that this is too strong a claim about geometric phases in
general: the underlying mechanism --- holonomy of a connection around a loop in
parameter or shape space, decoupled from the local dynamics along the way --- is
already fully present, and was fully worked out, in ordinary classical
mechanics. What is genuinely new in the quantum case is the specific bundle and
connection defined by the adiabatic theorem applied to a Hilbert space of
states; the geometric phenomenon of holonomy itself is not new.

\subsection{Open systems}

Two classical difficulties are worth setting side by side. The first is
statistical: Boltzmann's $H$-theorem derives macroscopic irreversibility from
microscopically time-reversible mechanics, and was met immediately by
Loschmidt's reversibility objection (reversing all velocities should reverse
the approach to equilibrium) and Zermelo's recurrence objection (Poincar\'e
recurrence guarantees a return arbitrarily close to any earlier state). Both
objections are correct as stated, and reconciling them with the observed arrow
of time is a genuinely hard, and still actively discussed, problem --- entirely
classical, with no reference to measurement or wavefunction collapse.

The second is dynamical: an accelerating point charge radiates, and Newton's
second law must be supplemented by the Abraham--Lorentz self-force,
\begin{equation}
  m\dot v = F_{\text{ext}} + \frac{2}{3}\frac{e^2}{c^3}\,\dot a \, ,
\end{equation}
a third-order correction proportional to the derivative of the acceleration.
This equation admits runaway solutions, in which the particle self-accelerates
exponentially with no external force at all; demanding that solutions remain
bounded, instead, forces the particle to begin accelerating \emph{before} the
external force is applied --- a pre-acceleration that appears to violate
causality, entirely within classical electrodynamics.

Both irreversibility and an apparent tension with causal, deterministic
evolution are frequently treated as problems that quantum theory introduces
--- through the measurement postulate, decoherence, or the collapse of the
wavefunction --- or that relativity introduces, through its reworking of
simultaneity. Loschmidt, Zermelo, and Abraham--Lorentz show that a purely
classical system, governed by Newton's laws and Maxwell's equations and nothing
else, already exhibits dissipative, apparently irreversible, and apparently
acausal behavior that has never been fully tamed.

\subsection{Looking Ahead}

Each of the four cases above is the subject of a full chapter later in this
book (chapters \ref{hm},\ref{nh},\ref{shape} and \ref{open})
, where the mathematics is developed in detail and the historical debate is
traced more carefully. The point of gathering them here, in outline, is
methodological: before a ``paradox'' encountered in relativity or in quantum
mechanics is attributed to relativity or to quantum mechanics as such, it is
worth checking whether classical mechanics, unaided, already contains a version
of it. Where it does, what is new in the generalization is not the paradox but,
at most, its dressing --- and the genuinely new content of the generalization
has to be located elsewhere.

\chapter{Bodies That Change Shape}
\label{shape}

\section{The falling cat}

Drop a cat, upside down, from rest. In free fall there is no
external torque about the center of mass, so angular momentum is
conserved --- and it starts, and remains, at zero. Naive intuition
says that a body with identically zero angular momentum cannot
rotate. And yet the cat lands on its feet, having reoriented itself
by roughly $180^\circ$. The naive inference silently assumes a
\emph{rigid} body; it simply does not apply to a body that can
deform itself internally, and the cat's righting maneuver --- a
sequence of twists of its front and back halves relative to one
another --- is a real, controlled deformation, not a violation of
any conservation law.

\section{Shape space and the gauge structure of deformation}

The clean way to see why this works, due to Shapere and
Wilczek~\cite{shapere+wilczek} and to Montgomery's later "gauge
theory of the falling cat"~\cite{montgomery}, is to separate a
deformable body's configuration into its \emph{shape} $s$ (the
internal, orientation-independent degrees of freedom --- relative
joint angles, say) and an overall orientation $R\in SO(3)$ that fixes
that shape in space. Locally, the full configuration space is a
principal fiber bundle with structure group $SO(3)$ over shape space,
and the requirement that total angular momentum stay identically
zero throughout the motion is a linear relation between the shape
velocity $\dot s$ and the angular velocity $\omega$ of the body,
\begin{equation}
  \omega \;=\; -\,\mathcal A(s)\,\dot s ,
  \label{connection}
\end{equation}
where $\mathcal A(s)$ is built from the shape-dependent inertia
tensor and its couplings. This is precisely the data of a connection
(a gauge potential) on the bundle: given any path $s(t)$ in shape
space, \eqref{connection} lifts it uniquely to a path in the full
configuration space, exactly as a connection lifts a path in the
base of a fiber bundle to a horizontal path upstairs.

The consequence that matters here is that if the shape returns
exactly to where it started after a closed loop in shape space,
$s(T) = s(0)$, the orientation need \emph{not} return to where it
started. The net rotation accumulated around the loop is the
\emph{holonomy} of the connection~\eqref{connection}, and for a
small loop enclosing area $\Sigma$ it is governed, Stokes'-theorem
style, by the curvature $\mathcal F = d\mathcal A + \mathcal
A\wedge\mathcal A$ of that connection:
\begin{equation}
  \Delta R \;\approx\; \exp\!\left[\oint_{\partial\Sigma}\mathcal
  A\right] \;\approx\; \exp\!\left[\int_{\Sigma}\mathcal F\right].
  \label{holonomy}
\end{equation}
This is exactly the structure of a non-Abelian geometric (Berry,
Wilczek--Zee) phase from quantum mechanics, and of Aharonov--Bohm
holonomy in gauge theory --- reproduced in full by an entirely
classical, entirely mechanical system. It is also, viewed from the
vantage point of Section~\ref{cm:nh}, nothing but another
instance of a nonholonomic constraint: "total angular momentum
$\equiv 0$" is a Pfaffian relation between $\dot s$ and $\omega$ that
fails to be integrable to any function of shape and orientation
alone \emph{exactly when} the connection~\eqref{connection} has
nonzero curvature. The falling-cat effect and the rolling coin are,
formally, the same phenomenon in different clothing.

\section{Swimming without pushing}

The same structure governs self-propulsion in a highly viscous
fluid, where inertia is negligible and net displacement over one
cycle of body deformation is possible only if the cycle of shapes is
genuinely non-reciprocal. This is Purcell's scallop theorem: a
scallop that opens and closes through the \emph{same} sequence of
shapes, forwards then backwards, achieves zero net displacement,
because time-reversing the (inertia-free) equations of motion
exactly reverses the trajectory. Net swimming is again a
holonomy/geometric-phase effect in an appropriately defined shape
bundle~\cite{shapere+wilczek-swim}.

Perhaps the most striking illustration of the underlying idea is Jack
Wisdom's result that a quasi-rigid body executing a cyclic sequence
of internal shape changes in a curved \emph{spacetime} --- as
prescribed by general relativity --- can achieve a net translation
even in complete vacuum, with nothing to push against at
all~\cite{wisdom}. The effect scales with the local spacetime
curvature and with the extent of the body's internal motion; it is
minuscule for any realistic body and is not a practical propulsion
mechanism, but as a matter of principle it is "swimming" through
spacetime using exactly the same holonomy mechanism as the falling
cat, now realized by the genuine nonlinearity of general relativity
rather than by a gauge potential on a shape bundle constructed by
hand. Once again, an outcome that sounds like it must violate
momentum conservation turns out to be perfectly consistent with
conservation laws, provided one tracks correctly what "conservation"
does and does not forbid for a body that is allowed to change shape.

\chapter{Open Systems}
\label{open}

\section{Variable mass and the Meshchersky equation}

Ordinary Newtonian mechanics implicitly assumes a fixed set of
particles. Rockets, raindrops growing by accretion, chains lifted
off a pile, and avalanches all violate that assumption: the mass $m$
of the body under study is itself a function of time. The correct
generalization of Newton's second law to such an \emph{open system}
--- worked out around the turn of the twentieth century by
I.\,V.\ Meshchersky --- is
\begin{equation}
  m\,\frac{d\mathbf v}{dt} \;=\; \mathbf F_{\mathrm{ext}}
  + \mathbf u\,\frac{dm}{dt} ,
  \label{meshchersky}
\end{equation}
where $\mathbf u$ is the velocity of the incoming or outgoing mass
\emph{relative to the body itself}, not relative to the lab frame.
Setting $\mathbf F_{\mathrm{ext}}=0$ and integrating
\eqref{meshchersky} with $\mathbf u$ constant and antiparallel to
the direction of travel yields the Tsiolkovsky rocket equation,
\begin{equation}
  v(t) - v(0) \;=\; u\,\ln\frac{m(0)}{m(t)} .
  \label{tsiolkovsky}
\end{equation}
The resemblance to Section~\ref{cm:hertz} is not superficial. The
term $\mathbf u\,\dot m$ in \eqref{meshchersky} looks, from the
point of view of the rocket body alone, exactly like an applied
force --- "thrust" --- and is routinely called one. But viewed as
part of the larger \emph{closed} system consisting of the rocket
\emph{together with} every parcel of exhaust it has ever ejected,
whose total momentum is exactly conserved, there is no force there
at all: only the ordinary transfer of momentum between two pieces of
one conserved whole. What looks like a force is, once again, an
artifact of insisting on an equation of motion for part of a system
rather than for the whole of it.

This is also a genuine pedagogical trap, worth flagging explicitly.
It is tempting to write $\mathbf F = d(m\mathbf v)/dt = m\dot{\mathbf
v} + \mathbf v\,\dot m$ and treat this directly as the equation of
motion of the open body. That expression is correct as a statement
about the rate of change of the momentum of the body \emph{alone},
but it is not, in general, equivalent to \eqref{meshchersky} and
gives the wrong dynamics unless the accreted or ejected mass happens
to leave (or join) with zero velocity relative to the body, i.e.\
$\mathbf u \equiv \mathbf v$ --- a special case, not the generic one,
and not the one realized by a rocket engine, whose whole purpose is
to eject mass at high relative velocity.

\section{Dissipative systems: putting friction back into a
  variational principle}

A second, quite different kind of "openness" arises whenever a
system exchanges energy with an environment that is not explicitly
modelled --- friction, drag, radiation reaction. None of these
forces derive from a potential in the ordinary sense, so the
Lagrangian and Hamiltonian formalisms, built for closed and
energy-conserving systems, do not apply directly. Two classical
devices repair this.

The first is Rayleigh's dissipation function $\mathcal
F_{\mathrm{diss}}(\dot q)$, quadratic in the generalized velocities
for linear damping, which modifies the Euler--Lagrange equations to
\begin{equation}
  \frac{d}{dt}\frac{\partial L}{\partial \dot q^i}
  - \frac{\partial L}{\partial q^i}
  \;=\; -\,\frac{\partial \mathcal F_{\mathrm{diss}}}{\partial \dot
  q^i}.
  \label{rayleigh}
\end{equation}
This correctly reproduces the (irreversible) equations of motion,
but it is a hybrid: $\mathcal F_{\mathrm{diss}}$ is not itself part
of a genuine, self-contained stationary-action principle for $L$
alone.

The second, due to Bateman (1931), is more radical and more
illuminating: a fully conservative, symplectic Lagrangian \emph{can}
be written for a damped oscillator, but only at the price of
doubling the number of degrees of freedom. Introduce an auxiliary
"mirror" coordinate $y$ alongside the physical coordinate $x$,
\begin{equation}
  L_{\mathrm{dual}}(x,\dot x,y,\dot y) \;=\;
  m\,\dot x\,\dot y \;+\; \frac{\gamma}{2}\bigl(x\dot y - \dot x
  y\bigr) \;-\; k\,xy ,
  \label{bateman}
\end{equation}
whose Euler--Lagrange equations are precisely the damped oscillator,
$\ddot x + (\gamma/m)\dot x + \omega_0^2 x = 0$, for the physical
coordinate $x$, paired with a time-reversed, exponentially
\emph{amplified} equation for $y$. The energy the
physical oscillator loses is exactly the energy its mirror partner
gains, so the doubled system is conservative even though either half
alone is not. This is a small, entirely classical and elementary
demonstration of an idea that later became central to the theory of
open quantum systems: dissipation and irreversibility in a subsystem
are not exceptions to an underlying conservative dynamics, but the
generic appearance of a conservative dynamics once part of the full
system --- here, quite literally, the mirror half --- has been
discarded from the description.

\section{When "open" turns pathological: radiation reaction}

Not every attempt to write a closed-looking equation for an open
system ends benignly. An accelerating point charge radiates, and by
momentum conservation it must experience a recoil, or self-force, in
consequence. Treating the charge \emph{in isolation} --- an open
system in the strict sense, since the charge together with its own
radiated field is really a system with infinitely many degrees of
freedom, formally "integrated out" down to a single effective force
term --- leads to the (nonrelativistic) Abraham--Lorentz equation,
\begin{equation}
  m\,\ddot{\mathbf x} \;=\; \mathbf F_{\mathrm{ext}} \;+\;
  m\tau\,\dddot{\mathbf x} ,
  \qquad
  \tau \;=\; \frac{q^2}{6\pi\varepsilon_0\,m\,c^3}
  \;=\; \frac{2}{3}\,\frac{r_e}{c},
  \label{abraham-lorentz}
\end{equation}
where $r_e$ is the classical electron radius. Because
\eqref{abraham-lorentz} is third order in position, it requires
more initial data than position and velocity to determine a unique
solution, and generic initial data produce \emph{runaway} solutions
in which the self-acceleration grows without bound even with
$\mathbf F_{\mathrm{ext}}=0$. Suppressing the runaway by hand forces
the opposite pathology: the charge must begin accelerating
\emph{before} an external force is switched on --- a small, short-lived,
but genuinely acausal-looking pre-acceleration. (In practice, both
pathologies signal that \eqref{abraham-lorentz} is being pushed
outside its domain of validity, on timescales comparable to $\tau
\sim 10^{-24}\,\mathrm s$; treated properly, as an effective equation
valid only for slowly varying external forces --- via, e.g., the
Landau--Lifshitz reduction of order --- the pathologies disappear.)
The lesson is a cautionary complement to the rest of this section:
"integrating out" the unmodelled part of an open system is not
always harmless bookkeeping. Sometimes it is, as with the Meshchersky
equation; sometimes, as here, it produces an equation that is simply
malformed unless handled with real care.

\chapter{Nonholonomic Constraints}
\label{nh}

\section{Holonomic versus nonholonomic}

A constraint of the form $f(q,t)=0$, where $q=(q^1,\dots,q^n)$ are
generalized coordinates, is called \emph{holonomic}: it can be
solved, at least locally, to eliminate one coordinate outright, and
it reduces the dimension of the configuration space itself. A more
general class of constraints is linear in the velocities,
\begin{equation}
  \sum_{i=1}^n a_i(q,t)\,\dot q^i + a_t(q,t) \;=\;0,
  \label{pfaffian}
\end{equation}
a so-called Pfaffian constraint. If the left-hand side of
\eqref{pfaffian} happens to be the total time derivative of some
function $f(q,t)$, the constraint is secretly holonomic in disguise.
If it is not --- if no such $f$ exists --- the constraint is called
\emph{nonholonomic}: it restricts the \emph{velocities} available at
each point of $Q$, without restricting which \emph{configurations}
are reachable.

The canonical example is a coin (or disk) of radius $R$ rolling
without slipping on a plane, with contact point $(x,y)$, heading
angle $\phi$, and roll angle $\theta$ as coordinates. The
rolling condition reads
\begin{equation}
  \dot x = R\,\dot\theta\cos\phi ,\qquad
  \dot y = R\,\dot\theta\sin\phi ,
  \label{eq:rolling-coin}
\end{equation}
which is not integrable: it constrains the instantaneous velocity to
a two-dimensional subspace of the four-dimensional tangent space at
each point, and yet, by alternating steering and rolling, the coin
can be driven from \emph{any} configuration $(x,y,\phi)$ to any
other. The instantaneous constraint is genuine, but it places no
restriction whatsoever on the reachable set of configurations --- an
instance of what control theorists call bracket-generating, or
Chow--Rashevski\u{i}, controllability. This already runs against a
piece of intuition carried over uncritically from the holonomic
case, where a nontrivial velocity constraint always shrinks the
accessible configuration space too.

A second standard example, the \emph{Chaplygin sleigh}, makes an
even sharper point. It is a rigid body sliding frictionlessly on a
plane, with a knife edge that permits motion only along its own
length (no sideways slip). The system conserves energy exactly ---
there is no friction anywhere in the model --- and yet its reduced
dynamics on the velocity variables is attracted to a fixed
equilibrium: left alone, the sleigh asymptotically settles into
straight-line motion with no rotation, knife edge trailing the
center of mass~\cite{chaplygin}. This
happens because the reduced equations of motion on the nonholonomic
constraint surface, unlike those of an ordinary holonomic mechanical
system, do not in general preserve phase-space (Liouville) volume,
even though they do preserve energy. A perfectly conservative
mechanical system can therefore display genuinely dissipative-looking
long-time behaviour --- attracting fixed points, basins of
attraction --- with no friction anywhere in the equations. Energy
conservation, by itself, is not enough to rule out the qualitative
signatures usually associated with dissipation once constraints stop
being holonomic.

\section{Two variational principles, two different theories}

For holonomic constraints, two routes to the equations of motion ---
D'Alembert's principle (constraint forces do no virtual work, for
virtual displacements compatible with the instantaneous constraint)
and Hamilton's principle (the true path extremizes the action among
\emph{all} paths satisfying the constraint) --- are provably
equivalent. For nonholonomic (non-integrable, velocity-linear)
constraints, this equivalence simply fails, and the failure is not a
minor technicality.

There are two natural-looking ways to generalize the variational
treatment to a system with a nonholonomic constraint of the
form~\eqref{pfaffian}.

\begin{description}
  \item[Nonholonomic (Lagrange--d'Alembert, or Chetaev) mechanics.]
  Impose the constraint directly on virtual displacements at each
  instant, exactly as in the holonomic case, and only afterwards
  derive equations of motion:
  \begin{equation}
    \frac{d}{dt}\frac{\partial L}{\partial \dot q^i}
    - \frac{\partial L}{\partial q^i}
    \;=\; \sum_a \lambda_a\, a_i^a(q),
    \qquad
    \sum_i a_i^a(q)\,\dot q^i + a_t^a(q) = 0 .
    \label{lagrange-dalembert}
  \end{equation}
  \item[Vakonomic mechanics.] (The name --- "variational axiomatic
  kind" --- was coined by Kozlov~\cite{kozlov}.)
  Instead, treat the nonholonomic constraint as a genuine constraint
  on the space of \emph{paths} from the outset, as one would an
  isoperimetric constraint in the calculus of variations, and only
  then extremize the action over that constrained set of paths.
\end{description}

For holonomic constraints these two prescriptions coincide. For
genuinely nonholonomic constraints they generically do not: it has
been shown, for instance, that for the rolling coin on an inclined
plane no nontrivial vakonomic solution reproduces the
Lagrange--d'Alembert solution at
all~\cite{zampieri}. This is not merely a
mathematical curiosity to be settled by taste: Lewis and Murray
built and instrumented a physical system --- a ball rolling without
slipping on a rotating table --- specifically to adjudicate between
the two theories, and found that the observed motion matches the
Lagrange--d'Alembert (nonholonomic) prediction; they were unable to
produce vakonomic simulations resembling the experimental data at
all~\cite{lewis+murray}. The variational principle that looks like
the most direct generalization of Hamilton's principle to velocity
constraints is, for constraints realized by an actual rolling or
sliding contact, simply the wrong physical theory.

This does not make vakonomic mechanics useless --- it correctly
describes optimal-control and sub-Riemannian-geodesic problems, where
the "constraint" genuinely is a restriction on the class of paths
being optimized over, rather than a constraint enforced
instant-by-instant by a physical reaction force. The moral is
subtler and more interesting than "one theory is right and the other
wrong": the bare constraint equation~\eqref{pfaffian}, by itself,
does not determine the dynamics. The correct equations of motion
depend on how the constraint is physically \emph{realized} ---
information that is invisible at the level of the constraint
equation alone, and that a first course in mechanics, built entirely
on holonomic examples where this ambiguity cannot arise, gives no
reason to expect.
